% GENERATED FILE. Edit canonical lessons, then run npm run generate:books.
% canonical-source-hash: fnv1a64:c59e49002b1eb1bf
% canonical-lessons: UR-C310-talvar, UR-C310-dhal, UR-C310-bhala, UR-C310-kaman, UR-C310-lashkar

\chapter{Swords and Shields}
\label{ch:ur-swords-and-shields}
\hlchaptermodality{310}

\begin{chapteropening}
\textbf{By the end of this chapter:} \emph{I can describe how battles were fought long ago.}

Complete the last lesson of chapter 310: I can describe how battles were fought long ago.
\end{chapteropening}

\section[\texorpdfstring{\ur{تلوار} (talvār)}{تلوار (talvār)}]{\ur{تلوار} (talvār) — a sword}

\label{lesson:UR-C310-talvar}



\begin{quote}
Before the new one: say the Urdu for a cannon, then the Urdu for to clash.
\end{quote}

\subsection*{You'll want to know: \ur{تلوار}}
\textbf{\ur{تلوار}} — \emph{talvār} — \textquotedblleft{}a sword\textquotedblright{}.

\begin{quote}
\ur{سپاہی} \ur{کے} \ur{پاس} \ur{تلوار} \ur{تھی۔}
\end{quote}

\begin{quote}
\emph{sipāhī ke pās talvār thī.}
\end{quote}

\begin{quote}
The soldier had a sword.
\end{quote}

Feminine.

\begin{cousinweb}
\textbf{\ur{تلوار}} (\emph{talvār}) is inherited from Sanskrit \emph{taravāri}, \textquotedblleft{}a sword\textquotedblright{} (CDIAL 5706), where CDIAL lists the Hindi cousin \emph{talwār}. It is feminine.
\end{cousinweb}

\begin{grammarlens}[title={What you've built}]
One more word for telling what happened.
\end{grammarlens}

\subsection*{Your turn}
\begin{itemize}
\raggedright
  \item \emph{Say it:} \emph{talvār}
  \item \emph{Say it:} \emph{talvār}, once more
  \item \emph{Say it:} \emph{sipāhī ke pās talvār thī.} — The soldier had a sword.
  \item \emph{Recall:} say the Urdu for a cannon, then the Urdu for to clash, then say \emph{talvār} again
\end{itemize}

\begin{tcolorbox}[breakable,colback=teal!4,colframe=teal!35!black,title={Before you move on}]
What does \ur{تلوار} mean? (\textquotedblleft{}A sword\textquotedblright{}.) Say it once more.
\end{tcolorbox}
\section[\texorpdfstring{\ur{ڈھال} (ḍhāl)}{ڈھال (ḍhāl)}]{\ur{ڈھال} (ḍhāl) — a shield}

\label{lesson:UR-C310-dhal}



\begin{quote}
Before the new one: say the Urdu for to clash, then the Urdu for a sword.
\end{quote}

\subsection*{You'll want to know: \ur{ڈھال}}
\textbf{\ur{ڈھال}} — \emph{ḍhāl} — \textquotedblleft{}a shield\textquotedblright{}.

\begin{quote}
\ur{اس} \ur{نے} \ur{ڈھال} \ur{اٹھائی۔}
\end{quote}

\begin{quote}
\emph{us ne ḍhāl uṭhāī.}
\end{quote}

\begin{quote}
He raised his shield.
\end{quote}

Feminine.

\begin{cousinweb}
\textbf{\ur{ڈھال}} (\emph{ḍhāl}) is inherited from Sanskrit \emph{ḍhāla}, \textquotedblleft{}a shield\textquotedblright{} (CDIAL 5583). It is feminine: \textbf{\ur{ڈھال} \ur{اٹھائی}}.
\end{cousinweb}

\begin{grammarlens}[title={What you've built}]
One more word for telling what happened.
\end{grammarlens}

\subsection*{Your turn}
\begin{itemize}
\raggedright
  \item \emph{Say it:} \emph{ḍhāl}
  \item \emph{Say it:} \emph{ḍhāl}, once more
  \item \emph{Say it:} \emph{us ne ḍhāl uṭhāī.} — He raised his shield.
  \item \emph{Recall:} say the Urdu for to clash, then the Urdu for a sword, then say \emph{ḍhāl} again
\end{itemize}

\begin{tcolorbox}[breakable,colback=teal!4,colframe=teal!35!black,title={Before you move on}]
What does \ur{ڈھال} mean? (\textquotedblleft{}A shield\textquotedblright{}.) Say it once more.
\end{tcolorbox}
\section[\texorpdfstring{\ur{بھالا} (bhālā)}{بھالا (bhālā)}]{\ur{بھالا} (bhālā) — a spear}

\label{lesson:UR-C310-bhala}



\begin{quote}
Before the new one: say the Urdu for a sword, then the Urdu for a shield.
\end{quote}

\subsection*{You'll want to know: \ur{بھالا}}
\textbf{\ur{بھالا}} — \emph{bhālā} — \textquotedblleft{}a spear\textquotedblright{}.

\begin{quote}
\ur{اس} \ur{کے} \ur{ہاتھ} \ur{میں} \ur{لمبا} \ur{بھالا} \ur{تھا۔}
\end{quote}

\begin{quote}
\emph{us ke hāth meṅ lambā bhālā thā.}
\end{quote}

\begin{quote}
He had a long spear in his hand.
\end{quote}

Masculine.

\begin{cousinweb}
\textbf{\ur{بھالا}} (\emph{bhālā}) is inherited: CDIAL lists the Hindi cousin \emph{bhālā}, \textquotedblleft{}a large spear\textquotedblright{}, under Sanskrit \emph{bhalla} (CDIAL 9409).
\end{cousinweb}

\begin{grammarlens}[title={What you've built}]
One more word for telling what happened.
\end{grammarlens}

\subsection*{Your turn}
\begin{itemize}
\raggedright
  \item \emph{Say it:} \emph{bhālā}
  \item \emph{Say it:} \emph{bhālā}, once more
  \item \emph{Say it:} \emph{us ke hāth meṅ lambā bhālā thā.} — He had a long spear in his hand.
  \item \emph{Recall:} say the Urdu for a sword, then the Urdu for a shield, then say \emph{bhālā} again
\end{itemize}

\begin{tcolorbox}[breakable,colback=teal!4,colframe=teal!35!black,title={Before you move on}]
What does \ur{بھالا} mean? (\textquotedblleft{}A spear\textquotedblright{}.) Say it once more.
\end{tcolorbox}
\section[\texorpdfstring{\ur{کمان} (kamān)}{کمان (kamān)}]{\ur{کمان} (kamān) — a bow for arrows}

\label{lesson:UR-C310-kaman}



\begin{quote}
Before the new one: say the Urdu for a shield, then the Urdu for a spear.
\end{quote}

\subsection*{You'll want to know: \ur{کمان}}
\textbf{\ur{کمان}} — \emph{kamān} — \textquotedblleft{}a bow for arrows\textquotedblright{}.

\begin{quote}
\ur{سپاہی} \ur{کمان} \ur{چلانا} \ur{جانتا} \ur{تھا۔}
\end{quote}

\begin{quote}
\emph{sipāhī kamān chalānā jāntā thā.}
\end{quote}

\begin{quote}
The soldier knew how to use a bow.
\end{quote}

Feminine.

\begin{cousinweb}
\textbf{\ur{کمان}} (\emph{kamān}) is usually traced to Persian \emph{kamān}, \textquotedblleft{}a bow\textquotedblright{}. It is feminine.
\end{cousinweb}

\begin{grammarlens}[title={What you've built}]
One more word for telling what happened.
\end{grammarlens}

\subsection*{Your turn}
\begin{itemize}
\raggedright
  \item \emph{Say it:} \emph{kamān}
  \item \emph{Say it:} \emph{kamān}, once more
  \item \emph{Say it:} \emph{sipāhī kamān chalānā jāntā thā.} — The soldier knew how to use a bow.
  \item \emph{Recall:} say the Urdu for a shield, then the Urdu for a spear, then say \emph{kamān} again
\end{itemize}

\begin{tcolorbox}[breakable,colback=teal!4,colframe=teal!35!black,title={Before you move on}]
What does \ur{کمان} mean? (\textquotedblleft{}A bow for arrows\textquotedblright{}.) Say it once more.
\end{tcolorbox}
\section[\texorpdfstring{\ur{لشکر} (lashkar)}{لشکر (lashkar)}]{\ur{لشکر} (lashkar) — a host; an army on the march}

\label{lesson:UR-C310-lashkar}



\begin{quote}
Before the new one: say the Urdu for a spear, then the Urdu for a bow for arrows.
\end{quote}

\subsection*{You'll want to know: \ur{لشکر}}
\textbf{\ur{لشکر}} — \emph{lashkar} — \textquotedblleft{}a host; an army on the march\textquotedblright{}.

\begin{quote}
\ur{لشکر} \ur{شام} \ur{کو} \ur{شہر} \ur{پہنچا۔}
\end{quote}

\begin{quote}
\emph{lashkar shām ko shahr pahuṅchā.}
\end{quote}

\begin{quote}
The army reached the city in the evening.
\end{quote}

Masculine.

\begin{cousinweb}
\textbf{\ur{لشکر}} (\emph{lashkar}) is usually traced to Persian \emph{lashkar}, \textquotedblleft{}an army, a camp\textquotedblright{}. It is the old word of the war stories; \textbf{\ur{فوج}} is the army of today.
\end{cousinweb}

\begin{grammarlens}[title={What you've built}]
One more word for telling what happened.
\end{grammarlens}

\subsection*{Your turn}
\begin{itemize}
\raggedright
  \item \emph{Say it:} \emph{lashkar}
  \item \emph{Say it:} \emph{lashkar}, once more
  \item \emph{Say it:} \emph{lashkar shām ko shahr pahuṅchā.} — The army reached the city in the evening.
  \item \emph{Recall:} say the Urdu for a spear, then the Urdu for a bow for arrows, then say \emph{lashkar} again
\end{itemize}

\begin{tcolorbox}[breakable,colback=teal!4,colframe=teal!35!black,title={Before you move on}]
What does \ur{لشکر} mean? (\textquotedblleft{}A host; an army on the march\textquotedblright{}.) Say it once more.
\end{tcolorbox}
